\subsection{Modulus stabilisation in the multiple-modulus framework --- arXiv:2310.10369}

\textbf{Authors:} Stephen F. King, Xin Wang

\paragraph{主な主張}
dilaton の K\"{a}hler potential への非摂動補正を含む複数モジュラス模型で、追加の物質場に依存せず、$\tau_i=i,\omega$ の組合せに dS 極小を実現できることを示した \cite{King:2023snq}。

\paragraph{新規性}
単一モジュラスの解析を三つのモジュラスへ拡張し、真空配置ごとの安定化条件を分類して、複数の modular symmetry を用いるニュートリノ模型との対応を示した。

\paragraph{位置付け}
複数モジュラスの真空選択と flavor 現象論を結び付け、固定点での dS 真空を構成するための条件を整理した研究である。

\paragraph{コメント（読書メモ）}
読書時点での最近の moduli stabilization 研究が、非常に詳しく整理されていると感じた。
